\question[3]
Vergleiche Tiefensuche und Breitensuche.
\begin{parts}
\part Welche Datenstruktur nutzt die Tiefensuche, um sich die noch offenen Teilbäume zu merken, welche die Breitensuche?
\part Warum eignet sich für die Breitensuche keine Datenstruktur, die nach dem LIFO-Prinzip arbeitet?
\part Ein gesuchter Wert liegt nah an der Wurzel in einem sehr großen Baum. Welches Verfahren findet ihn in der Regel schneller? Begründe.
\end{parts}
\begin{solutionbox}{3.5cm}
a) Die Tiefensuche nutzt einen Keller, die Breitensuche eine Schlange. \\
b) Die Knoten sollen ebenenweise besucht werden. Wer zuerst eingereiht wurde, muss deshalb auch zuerst drankommen (FIFO). \\
c) Die Breitensuche, denn sie besucht die Knoten Ebene für Ebene. Die Tiefensuche läuft vielleicht erst weit in einen falschen Teilbaum hinein.
\end{solutionbox}
